\section{Encoding-Regime Separation}\label{sec:dichotomy}

Formally: this section proves regime-indexed access separation for a fixed sufficiency predicate, then lifts bit-operation bounds to thermodynamic/accounting consequences.

The hardness results of Section~\ref{sec:hardness} apply to worst-case instances under the succinct encoding. This section states an encoding-regime separation: an explicit-state upper bound versus a succinct-encoding worst-case lower bound, and an intermediate query-access obstruction family. The complexity class is a property of the encoding measured against a fixed decision question (Remark~\ref{rem:question-vs-problem}), not a property of the question itself: the same sufficiency question admits polynomial-time algorithms under one encoding and worst-case exponential cost under another.
Within the paper's full regime map, this section is the static access-layer decomposition; theorem-level stochastic and sequential core placements are given in Sections~\ref{sec:stochastic} and~\ref{sec:sequential}, with strict inter-regime separation in Section~\ref{sec:regime-hierarchy}.

Physical reading: the separation is an access-and-budget separation for the same decision semantics. Explicit access assumes full boundary observability; succinct access assumes compressed structure that may require exponential expansion to certify universally; query access assumes bounded measurement interfaces, and the obstruction persists across value-entry/state-batch interfaces with oracle-lattice transfer/strictness closures.\claimstamp{T\ref{thm:dichotomy};P\ref{prop:query-regime-obstruction},\ref{prop:query-value-entry-lb},\ref{prop:query-state-batch-lb},\ref{prop:oracle-lattice-transfer},\ref{prop:oracle-lattice-strict},\ref{prop:physical-claim-transport}}{E,S+ETH,Qf,Qb,AR}
\leanmeta{\LH{CC16}, \LH{CC50}, \LHrng{L}{8}{9}, \LHrng{L}{20}{21}, \LH{CT6}.}

\paragraph{Model note.} Theorem~\ref{thm:dichotomy} has Part~1 in [E] (time polynomial in $|S|$) and Part~2 in [S+ETH] (time exponential in $n$). Proposition~\ref{prop:query-regime-obstruction} is [Q\_fin] (finite-state Opt-oracle core), while Propositions~\ref{prop:query-value-entry-lb} and~\ref{prop:query-state-batch-lb} are [Q\_bool] interface refinements (value-entry and state-batch). Relative to Definition~\ref{def:regime-simulation}, Propositions~\ref{prop:query-subproblem-transfer} and~\ref{prop:oracle-lattice-transfer} are explicit simulation-transfer instances. The [E] vs [S+ETH] split in Theorem~\ref{thm:dichotomy} is a same-question encoding/cost-model separation, not a bidirectional simulation claim between those regimes. These regimes are defined in Section~\ref{sec:encoding} and are not directly comparable as functions of one numeric input-length parameter.

\begin{theorem}[Explicit--Succinct Regime Separation]\label{thm:dichotomy}
Let $\mathcal{D} = (A, X_1, \ldots, X_n, U)$ be a decision problem with $|S| = N$ states. Let $k^*$ be the size of the minimal sufficient set.

\begin{enumerate}
\item \textbf{[E] Explicit-state upper bound:} Under the explicit-state encoding, SUFFICIENCY-CHECK is solvable in time polynomial in $N$ (e.g. $O(N^2|A|)$).

\item \textbf{[S+ETH] Succinct lower bound (worst case):} Assuming ETH, there exists a family of succinctly encoded instances with $n$ coordinates and minimal sufficient set size $k^* = n$ such that SUFFICIENCY-CHECK requires time $2^{\Omega(n)}$.
\end{enumerate}
\leanmeta{\LH{CC16}, \LHrng{CC}{23}{24}.}
\end{theorem}

\begin{proof}
\textbf{Part 1 (Explicit-state upper bound):} Under the explicit-state encoding, SUFFICIENCY-CHECK is decidable in time polynomial in $N$ by direct enumeration: compute $\Opt(s)$ for all $s\in S$ and then check all pairs $(s,s')$ with $s_I=s'_I$.

Equivalently, for any fixed $I$, the projection map $s\mapsto s_I$ has image size $|S_I|\le |S|=N$, so any algorithm that iterates over projection classes (or over all state pairs) runs in polynomial time in $N$. Thus, in particular, when $k^*=O(\log N)$, SUFFICIENCY-CHECK is solvable in polynomial time under the explicit-state encoding.

\paragraph{Remark (bounded coordinate domains).} In the general model $S=\prod_i X_i$, for a fixed $I$ one always has $|S_I|\le \prod_{i\in I}|X_i|$ (and $|S_I|\le N$). If the coordinate domains are uniformly bounded, $|X_i|\le d$ for all $i$, then $|S_I|\le d^{|I|}$.

\textbf{Part 2 (Succinct ETH lower bound, worst case):} A strengthened version of the TAUTOLOGY reduction used in Theorem~\ref{thm:sufficiency-conp} produces a family of instances in which the minimal sufficient set has size $k^* = n$: given a Boolean formula $\varphi$ over $n$ variables, we construct a decision problem with $n$ coordinates such that if $\varphi$ is not a tautology then \emph{every} coordinate is decision-relevant (thus $k^*=n$). This strengthening is mechanized in Lean (see Appendix~\ref{app:lean}). Under the Exponential Time Hypothesis (ETH) \cite{impagliazzo2001complexity,arora2009computational}, TAUTOLOGY requires time $2^{\Omega(n)}$ in the succinct encoding, so SUFFICIENCY-CHECK inherits a $2^{\Omega(n)}$ worst-case lower bound via this reduction.
\end{proof}

\begin{corollary}[Regime Separation (by Encoding)]
There is a clean separation between explicit-state tractability and succinct worst-case hardness (with respect to the encodings in Section~\ref{sec:encoding}):
\begin{itemize}
\item Under the explicit-state encoding, SUFFICIENCY-CHECK is polynomial in $N=|S|$.
\item Under ETH, there exist succinctly encoded instances with $k^* = \Omega(n)$ (indeed $k^*=n$) for which SUFFICIENCY-CHECK requires $2^{\Omega(n)}$ time.
\end{itemize}
For Boolean coordinate spaces ($N = 2^n$), the explicit-state bound is polynomial in $2^n$ (exponential in $n$), while under ETH the succinct lower bound yields $2^{\Omega(n)}$ time for the hard family in which $k^* = \Omega(n)$.
\end{corollary}

\begin{proposition}[Finite-State Query Lower-Bound Core via Empty-Set Subproblem]\label{prop:query-regime-obstruction}
In the finite-state query regime [Q\_fin], let $S$ be any finite state type with $|S|\ge 2$ and let $Q\subset S$ with $|Q|<|S|$.
Then there exist two decision problems $\mathcal{D}_{\mathrm{yes}},\mathcal{D}_{\mathrm{no}}$ over the same state space such that:
\begin{itemize}
\item their oracle views on all states in $Q$ are identical,
\item $\emptyset$ is sufficient for $\mathcal{D}_{\mathrm{yes}}$,
\item $\emptyset$ is not sufficient for $\mathcal{D}_{\mathrm{no}}$.
\end{itemize}
Consequently, no deterministic query procedure using fewer than $|S|$ state queries can solve SUFFICIENCY-CHECK on all such instances; the worst-case query complexity is $\Omega(|S|)$.
\leanmeta{\LH{CC50}, \LH{CC50}.}
\end{proposition}

\begin{proof}
This is exactly the finite-state indistinguishability theorem
\LH{CC4}. For any $|Q|<|S|$, there is an unqueried hidden state $s_0$ producing oracle-indistinguishable yes/no instances with opposite truth values on the $I=\emptyset$ subproblem. Since SUFFICIENCY-CHECK contains that subproblem, fewer than $|S|$ queries cannot be correct on both instances, yielding the $\Omega(|S|)$ worst-case lower bound.
\end{proof}

\begin{proposition}[Checking--Witnessing Duality on the Empty-Set Core]\label{prop:checking-witnessing-duality}
\claimstamp{P\ref{prop:checking-witnessing-duality}}{H=query-lb}
Let $n\ge 1$ and let $P$ be a finite family of Boolean-state pairs.
Assume $P$ is sound for refuting empty-set sufficiency in the sense that for every
\(\Opt : \{0,1\}^n \to \{0,1\}\), whenever \(\emptyset\) is not sufficient, some pair in \(P\)
witnesses disagreement of \(\Opt\).
Then
\[
|P| \ge 2^{n-1}.
\]
Equivalently, any sound checker must inspect at least the witness budget
\(W(n)=2^{n-1}\), so checking budget \(T(n)\) satisfies \(T(n)\ge W(n)\).
\leanmeta{\LHrng{WD}{3}{5}.}
\end{proposition}

\begin{corollary}[Information Barrier in Query Regimes]\label{cor:information-barrier-query}
\claimstamp{C\ref{cor:information-barrier-query}}{H=query-lb}
For the fixed sufficiency question, finite query interfaces induce an information barrier:
\begin{itemize}
\item [Q\_fin] (Opt-oracle): with $|Q|<|S|$, there exist yes/no instances indistinguishable on all queried states but with opposite truth values on the $I=\emptyset$ subproblem.
\item [Q\_bool] value-entry and state-batch interfaces preserve the same obstruction at cardinality $<2^n$.
\end{itemize}
Hence the barrier is representational-access based: without querying enough of the hidden state support, exact certification is blocked by indistinguishability rather than by search-strategy choice.
\leanmeta{\LHrng{CC}{27}{29}.}
\end{corollary}

\begin{corollary}[Boolean-Coordinate Instantiation]\label{cor:query-obstruction-bool}
\claimstamp{C\ref{cor:query-obstruction-bool}}{H=query-lb}
In the Boolean-coordinate state space $S=\{0,1\}^n$, Proposition~\ref{prop:query-regime-obstruction} yields the familiar $\Omega(2^n)$ worst-case query lower bound for Opt-oracle access.
\leanmeta{\LHrng{CC}{49}{50}.}
\end{corollary}

\begin{proposition}[Value-Entry Query Lower Bound]\label{prop:query-value-entry-lb}
In the mechanized Boolean value-entry query submodel [Q\_bool], for any deterministic procedure using fewer than $2^n$ value-entry queries $(a,s)\mapsto U(a,s)$, there exist two queried-value-indistinguishable instances with opposite truth values for SUFFICIENCY-CHECK on the $I=\emptyset$ subproblem. Therefore worst-case value-entry query complexity is also $\Omega(2^n)$.
\leanmeta{\LH{CC29}, \LH{CC50}.}
\end{proposition}

\begin{proof}
The theorem \LH{CC50} constructs, for any query set of cardinality $<2^n$, an unqueried hidden state $s_0$ and a yes/no instance pair that agree on all queried values but disagree on $\emptyset$-sufficiency. The auxiliary bound \LH{CC50} ensures that fewer than $2^n$ value-entry queries cannot cover all states, so the indistinguishability argument applies in the worst case.
\end{proof}

\begin{proposition}[Subproblem-to-Full Transfer Rule]\label{prop:query-subproblem-transfer}
If every full-problem solver induces a solver for a fixed subproblem, then any lower bound for that subproblem lifts to the full problem.
\leanmeta{\LH{CC52}, \LHrng{CC}{58}{59}.}
\end{proposition}

This is the restriction-map instance of Definition~\ref{def:regime-simulation}: a solver for the full regime induces one for the restricted subproblem regime, so lower bounds transfer.

\begin{proposition}[Randomized Robustness (Seedwise)]\label{prop:query-randomized-robustness}
\claimstamp{P\ref{prop:query-randomized-robustness}}{H=query-lb}
In [Q\_bool], for any query set with cardinality $<2^n$, the indistinguishable yes/no pair from Proposition~\ref{prop:query-regime-obstruction} forces one decoding error \emph{per random seed} for any seed-indexed decoder from oracle transcripts.
Consequently, finite-support randomization does not remove the obstruction: averaging preserves a constant error floor on the hard pair.
\leanmeta{\LH{CC50}, \LH{CC50}.}
\end{proposition}

\begin{proposition}[Randomized Robustness (Weighted Form)]\label{prop:query-randomized-weighted}
\claimstamp{P\ref{prop:query-randomized-weighted}}{H=query-lb}
For any finite-support seed weighting $\mu$, the same hard pair satisfies a weighted identity:
the weighted sum of yes-error and no-error equals total seed weight. Hence randomization cannot collapse both errors simultaneously.
\leanmeta{\LH{CC50}, \LH{CC50}.}
\end{proposition}

\begin{proposition}[State-Batch Query Lower Bound]\label{prop:query-state-batch-lb}
In [Q\_bool], the same $\Omega(2^n)$ lower bound holds for a state-batch oracle that returns the full Boolean-action utility tuple at each queried state.
\leanmeta{\LH{CC28}, \LH{CC50}.}
\end{proposition}

\begin{proposition}[Finite-State Empty-Subproblem Generalization]\label{prop:query-finite-state-generalization}
\claimstamp{P\ref{prop:query-finite-state-generalization}}{H=query-lb}
The empty-subproblem indistinguishability lower-bound core extends from Boolean-vector state spaces to any finite state type with at least two states.
\leanmeta{\LH{CC50}, \LH{CC50}.}
\end{proposition}

\begin{proposition}[Adversary-Family Tightness by Full Scan]\label{prop:query-tightness-full-scan}
\claimstamp{P\ref{prop:query-tightness-full-scan}}{H=query-lb}
For the const/spike adversary family used in the query lower bounds, querying all states distinguishes the pair; thus the lower-bound family is tight up to constant factors under full-state scan.
\leanmeta{\LH{CC50}.}
\end{proposition}

\begin{proposition}[Weighted Query-Cost Transfer]\label{prop:query-weighted-transfer}
\claimstamp{P\ref{prop:query-weighted-transfer}}{H=query-lb}
Let $w(q)$ be per-query cost and $w_{\min}$ a lower bound on all queried costs.
Any cardinality lower bound $|Q|\geq L$ lifts to weighted cost:
\[
\sum_{q\in Q} w(q)\ \ge\ w_{\min}\cdot L.
\]
\leanmeta{\LH{CC50}, \LH{CC50}.}
\end{proposition}

\begin{proposition}[Oracle-Lattice Transfer: Batch $\Rightarrow$ Entry]\label{prop:oracle-lattice-transfer}
In [Q\_bool], agreement on state-batch oracle views over touched states implies agreement on all value-entry views for the corresponding entry-query set.
\leanmeta{\LH{CC43}, \LH{CC52}.}
\end{proposition}

This is the oracle-transducer instance of Definition~\ref{def:regime-simulation}: batch transcripts induce entry transcripts while preserving indistinguishability for the target sufficiency question.

\begin{proposition}[Oracle-Lattice Strictness: Opt vs Value Entries]\label{prop:oracle-lattice-strict}
`Opt`-oracle views are strictly coarser than value-entry views: there exist instances with identical `Opt` views but distinguishable value entries.
\leanmeta{\LH{CC27}, \LH{CC29}.}
\end{proposition}

\begin{remark}[Instantiation of Definitions~\ref{def:structural-complexity} and~\ref{def:representational-hardness}]
Theorem~\ref{thm:dichotomy} and Propositions~\ref{prop:query-regime-obstruction}--\ref{prop:query-value-entry-lb} keep the structural problem fixed (same sufficiency relation) and separate representational hardness by access regime: explicit-state access exposes the boundary $s \mapsto \Opt(s)$, finite-state query access already yields $\Omega(|S|)$ lower bounds at the Opt-oracle level (Boolean instantiation: $\Omega(2^n)$), value-entry/state-batch interfaces preserve the obstruction in [Q\_bool], and succinct access can hide structure enough to force ETH-level worst-case cost on a hard family.
\end{remark}

This regime-typed landscape identifies tractability under [E], a finite-state query lower-bound core under [Q\_fin] with Boolean interface refinements under [Q\_bool], and worst-case intractability under [S+ETH] for the same underlying decision question.

The structural reason for this separation is enumerable access. Under the explicit-state encoding, the mapping $s \mapsto \Opt(s)$ is directly inspectable and sufficiency verification reduces to scanning state pairs, at cost polynomial in $|S|$. Under succinct encoding, the circuit compresses an exponential state space into polynomial size; universal verification over that compressed domain, namely ``for all $s, s'$ with $s_I = s'_I$, does $\Opt(s) = \Opt(s')$?'', carries worst-case cost exponential in $n$, because the domain the circuit compressed away cannot be reconstructed without undoing the compression.

Informally: if you can see enough of the decision boundary, checking can be fast; with limited access, it becomes expensive.

\subsection{Budgeted Physical Crossover}
Formally: this subsection states budget-indexed representational feasibility splits and links them to integrity-preserving policy closure.

The [E] vs [S] split can be typed against an explicit physical budget. Let $E(n)$ and $S(n)$ denote explicit and succinct representation sizes for the same decision question at scale parameter $n$, and let $B$ be a hard representation budget.
These crossover claims are interpreted through the physical-to-core encoding discipline of Section~\ref{sec:physical-transport}: physical assumptions are transferred explicitly into core assumptions, and encoded physical counterexamples map to core failures on the encoded slice.\claimstamp{P\ref{prop:physical-claim-transport};C\ref{cor:physical-counterexample-core-failure}}{AR}
\leanmeta{\LH{CT4}, \LH{CT6}, \LH{CT8}.}

\begin{proposition}[Budgeted Crossover Witness]\label{prop:budgeted-crossover}
\claimstamp{P\ref{prop:budgeted-crossover}}{H=exp-lb-conditional}
If $E(n) > B$ and $S(n) \le B$ for some $n$, then explicit representation is infeasible at $(B,n)$ while succinct representation remains feasible at $(B,n)$.
\leanmeta{\LH{CC45}, \LH{PBC1}.}
\end{proposition}

\begin{proof}
This is exactly the definitional split encoded in \LH{PBC1}: explicit infeasibility is $B < E(n)$ and succinct feasibility is $S(n)\le B$. The theorem \LH{CC2} returns both conjuncts.
\end{proof}

\begin{proposition}[Above-Cap Crossover Existence]\label{prop:crossover-above-cap}
\claimstamp{P\ref{prop:crossover-above-cap}}{H=exp-lb-conditional}
Assume there is a global succinct-size cap $C$ and explicit size is unbounded:
\[
\forall n,\ S(n)\le C,\qquad
\forall B,\ \exists n,\ B < E(n).
\]
Then for every budget $B\ge C$, there exists a crossover witness at budget $B$.
\leanmeta{\LH{CC44}, \LH{PBC2}.}
\end{proposition}

\begin{proof}
Fix $B\ge C$. By unboundedness of $E$, choose $n$ with $B < E(n)$. By the cap assumption, $S(n)\le C\le B$, so $(B,n)$ is a crossover witness. The mechanized closure theorem is \LH{CC1}.
\end{proof}

\begin{proposition}[Crossover Does Not Imply Exact-Certification Competence]\label{prop:crossover-not-certification}
\claimstamp{P\ref{prop:crossover-not-certification}}{H=exp-lb-conditional}
Fix any collapse schema ``exact certification competence $\Rightarrow$ complexity collapse'' for the target exact-certification predicate. Under the corresponding non-collapse assumption, exact-certification competence is impossible even when a budgeted crossover witness exists.
\leanmeta{\LH{CC46}.}
\end{proposition}

\begin{proof}
The mechanized theorem \LH{CC3} bundles the crossover split with the same logical closure used in Proposition~\ref{prop:integrity-resource-bound}: representational feasibility and certification hardness are distinct predicates.
\end{proof}

The same closure form is isolated as a standalone physical-collapse chain: requirement-indexed collapse schemas (\LH{PH11}) induce no physical collapse at that requirement (\LH{PH12}); canonical specialization yields \(P{=}NP\) physical impossibility under the declared collapse map (\LH{PH13}, \LH{PH14}, \LH{PH15}). The collapse map itself is derived in an explicit SAT bridge (polytime SAT witness + reduction bridge + runtime-to-energy transfer) (\LH{PH16}, \LH{PH17}, \LH{PH18}, \LH{PH19}, \LH{PH20}, \LH{PH21}, \LH{PH22}), and the assumption boundary is explicit: under bounded budget, positive per-bit cost, and exponential required-op lower bound, claimed collapse forces at least one assumption failure (\LH{PH26}, \LH{PH27}); finite-budget necessity is witnessed by explicit unbounded-budget collapse countermodels (\LH{PH31}, \LH{PH32}, \LH{PH33}).

\begin{proposition}[Crossover Policy Closure Under Integrity]\label{prop:crossover-policy}
\claimstamp{P\ref{prop:crossover-policy}}{H=exp-lb-conditional}
In the certifying-solver model (Definitions~\ref{def:solver-integrity}--\ref{def:competence-regime}), assume:
\begin{itemize}
\item a budgeted crossover witness at $(B,n)$,
\item the same collapse implication as Proposition~\ref{prop:integrity-resource-bound},
\item a solver that is integrity-preserving for the declared relation.
\end{itemize}
Then either full in-scope non-abstaining coverage fails, or declared budget compliance fails.
Equivalently, in uncertified hardness-blocked regions, integrity is compatible with $\mathsf{ABSTAIN}/\mathsf{UNKNOWN}$ but not with unconditional exact-competence claims.
\leanmeta{\LH{CC47}.}
\end{proposition}

\begin{corollary}[Finite Budget Implies Eventual Exact-Certification Failure]\label{cor:finite-budget-threshold-impossibility}
\claimstamp{C\ref{cor:finite-budget-threshold-impossibility}}{H=exp-lb-conditional}
Fix a physical computer model with finite total budget and strictly positive per-bit physical cost.
If required operation count for exact certification has an exponential lower bound, then there exists a threshold \(n_0\) such that for all \(n \ge n_0\), required physical work exceeds budget.
For the canonical requirement model used in this paper, the same eventual-budget-failure conclusion holds.
\leanmeta{\LHrng{PBC}{7}{8}.}
\end{corollary}

\begin{proposition}[Least Divergence Point at Fixed Budget]\label{prop:least-divergence-point}
\claimstamp{P\ref{prop:least-divergence-point}}{H=exp-lb-conditional}
Fix a budgeted encoding model and budget \(B\). If a crossover witness exists at budget \(B\), then there is a least index \(n_{\mathrm{crit}}\) such that crossover holds at \(n_{\mathrm{crit}}\), and no smaller index has crossover.
\leanmeta{\LH{PBC5}.}
\end{proposition}

\begin{proposition}[Eventual Explicit Infeasibility from Monotone Growth]\label{prop:eventual-explicit-infeasibility}
\claimstamp{P\ref{prop:eventual-explicit-infeasibility}}{H=exp-lb-conditional}
If explicit size is monotone in \(n\) and already exceeds budget \(B\) at some \(n_0\), then explicit infeasibility holds for all \(n\ge n_0\).
\leanmeta{\LH{PBC7}.}
\end{proposition}

\begin{proposition}[Payoff Threshold as Eventual Split]\label{prop:payoff-threshold}
\claimstamp{P\ref{prop:payoff-threshold}}{H=exp-lb-conditional}
If, beyond \(n_0\), explicit encoding is infeasible while succinct encoding remains feasible, then every \(n\ge n_0\) is a crossover point.
\leanmeta{\LHrng{PBC}{8}{9}.}
\end{proposition}

\begin{corollary}[No Universal Survivor Without Succinct Bound]\label{cor:no-universal-survivor-no-succinct-bound}
\claimstamp{C\ref{cor:no-universal-survivor-no-succinct-bound}}{H=exp-lb-conditional}
If both explicit and succinct encodings are unbounded, then for every fixed budget \(B\), there exists an index where explicit is infeasible and an index where succinct is infeasible.
\leanmeta{\LH{PBC10}.}
\end{corollary}

\begin{proposition}[Policy Closure at and Beyond Divergence]\label{prop:policy-closure-beyond-divergence}
\claimstamp{P\ref{prop:policy-closure-beyond-divergence}}{H=exp-lb-conditional}
Under the same integrity-resource collapse assumption as Proposition~\ref{prop:crossover-policy}, policy closure holds at any fixed divergence point and uniformly beyond any threshold where crossover holds.
\leanmeta{\LHrng{PBC}{11}{12}.}
\end{proposition}

Informally: representational crossover can happen even when exact certification is blocked; under integrity this means abstention or reduced coverage.

\subsection{Thermodynamic Lift and Cost Accounting}
Formally: this subsection converts bit-operation lower bounds into energy/carbon lower bounds via linear conversion and additive accounting.

The crossover and hardness statements are complexity-theoretic.
The thermodynamic lift follows the same assumption discipline as the rest of the paper:
claims are exactly those derivable from the conversion model and bit-operation lower-bound premises.

\begin{proposition}[Bit-to-Energy/Carbon Lift]\label{prop:thermo-lift}
\claimstamp{P\ref{prop:thermo-lift}}{H=cost-growth}
Fix a thermodynamic conversion model with constants $(\lambda,\rho)$ where $\lambda$ lower-bounds energy per irreversible bit operation and $\rho$ lower-bounds carbon per joule. If a bit-operation lower bound $b_{\mathrm{LB}} \le b_{\mathrm{used}}$ holds, then:
\[
E_{\mathrm{LB}}=\lambda b_{\mathrm{LB}} \le \lambda b_{\mathrm{used}},\qquad
C_{\mathrm{LB}}=\rho E_{\mathrm{LB}} \le \rho(\lambda b_{\mathrm{used}}).
\]
\leanmeta{\LHrng{CC}{65}{66}.}
\end{proposition}

\begin{proposition}[Hardness--Thermodynamics Bundle]\label{prop:thermo-hardness-bundle}
\claimstamp{P\ref{prop:thermo-hardness-bundle}}{H=cost-growth}
Under the same non-collapse/collapse schema used by Proposition~\ref{prop:integrity-resource-bound}, exact-certification competence is impossible; combined with a declared bit-operation lower bound, this yields simultaneous lower bounds on energy and carbon in the declared conversion model.
\leanmeta{\LH{CC67}.}
\end{proposition}

\begin{proposition}[Mandatory Physical Cost]\label{prop:thermo-mandatory-cost}
\claimstamp{P\ref{prop:thermo-mandatory-cost}}{H=cost-growth}
Within the linear thermodynamic model, if
\[
\lambda > 0,\quad \rho > 0,\quad b_{\mathrm{LB}} > 0,
\]
then both lower bounds are strictly positive:
\[
E_{\mathrm{LB}}=\lambda b_{\mathrm{LB}} > 0,\qquad
C_{\mathrm{LB}}=\rho E_{\mathrm{LB}} > 0.
\]
\leanmeta{\LH{CC68}.}
\end{proposition}

\begin{proposition}[Conserved Additive Accounting]\label{prop:thermo-conservation-additive}
\claimstamp{P\ref{prop:thermo-conservation-additive}}{H=cost-growth}
Within the same linear model, lower-bound accounting is additive over composed bit-operation lower bounds:
\[
E_{\mathrm{LB}}(b_1+b_2)=E_{\mathrm{LB}}(b_1)+E_{\mathrm{LB}}(b_2),\qquad
C_{\mathrm{LB}}(b_1+b_2)=C_{\mathrm{LB}}(b_1)+C_{\mathrm{LB}}(b_2).
\]
\leanmeta{\LH{CC64}.}
\end{proposition}

\begin{corollary}[Neukart--Vinokur Thermodynamic Duality]\label{cor:neukart-vinokur}
\claimstamp{C\ref{cor:neukart-vinokur}}{AR}
Let $C$ denote a complexity coordinate with $C = b_{\mathrm{LB}}$ from Proposition~\ref{prop:query-regime-obstruction}. For $\lambda > 0$ in the conversion model, the thermodynamic duality
\[
dU \ge \lambda\,dC
\]
holds as a direct specialization of Proposition~\ref{prop:thermo-lift}.
\leanmeta{\LH{CC50}, \LH{CC65}.}
\end{corollary}

\begin{proof}
Set $C := b_{\mathrm{LB}} = \Omega(2^n)$ from the query obstruction (Proposition~\ref{prop:query-regime-obstruction}). By Proposition~\ref{prop:thermo-lift}, $E_{\mathrm{LB}} = \lambda\,b_{\mathrm{LB}}$. Substituting $C$ for $b_{\mathrm{LB}}$ yields $dU \ge dE_{\mathrm{LB}} = \lambda\,dC$. The inequality is strict when $\lambda > 0$ and $dC > 0$, which holds for [S+ETH] hard families with $k^* = n$ (Theorem~\ref{thm:dichotomy}).
\end{proof}

\subsubsection{Physical Grounding: Discrete State to Thermodynamic Cost}\label{sec:discrete-spacetime-chain}

The thermodynamic lift (Proposition~\ref{prop:thermo-lift}) rests on a derivation chain that traces computational discreteness through established physics to thermodynamic bounds. Each step is either proven or cited as accepted physics; the composition yields non-trivial reach.

\begin{proposition}[Discrete State Implies Discrete Effective Time]\label{prop:discrete-state-time}
\claimstamp{P\ref{prop:discrete-state-time}}{AR}
Let $\mathcal{S}$ be a finite state space with deterministic dynamics $\tau : \mathcal{S} \to \mathcal{S}$. If $\tau$ is non-trivial (i.e., $\exists s.\, \tau(s) \neq s$), then for any trajectory $\{s_t\}_{t \in \mathbb{N}}$, the set of transition points $T = \{t : s_{t+1} \neq s_t\}$ is countable.
\leanmeta{\LHrng{DS}{1}{2}.}
\end{proposition}

\begin{proof}
Immediate: $T \subseteq \mathbb{N}$ and $\mathbb{N}$ is countable. The real dynamics lives on the transition set, which is necessarily discrete.
\end{proof}

\begin{proposition}[Lorentz Invariance Forces Compatible Discreteness]\label{prop:lorentz-discrete}
\claimstamp{P\ref{prop:lorentz-discrete}}{AR}
Under Lorentz invariance~\cite{einstein1905electrodynamics,minkowski1908space}, if temporal intervals are quantized with minimal unit $\Delta t > 0$, then spatial intervals must be quantized with compatible unit $\Delta x = c\Delta t$. Conversely, discrete space implies discrete time.
\leanmeta{\LHrng{DS}{3}{4}.}
\end{proposition}

\begin{proof}
The spacetime interval $s^2 = c^2\Delta t^2 - \Delta x^2$ is Lorentz-invariant. If $\Delta x^2$ takes discrete values and $s^2$ is frame-independent, then $c^2\Delta t^2$ must also take discrete values. This is the core of Snyder's quantized spacetime~\cite{snyder1947quantized}: Lorentz-invariant discrete spacetime requires non-commutative coordinates with $[x_\mu, x_\nu] = i\ell^2 J_{\mu\nu}$, where $J_{\mu\nu}$ are Lorentz generators, enforcing compatible space-time discreteness.
\end{proof}

The Planck scale $\ell_P = \sqrt{\hbar G / c^3}$ is the unique minimal length derivable from dimensional analysis of the fundamental constants $\hbar$, $G$, and $c$~\cite{planck1899irreversible}.\claimstamp{}{AR} This is not proven but cited as established physics: no other combination of these constants yields a length.

\begin{proposition}[Computational-Thermodynamic Chain]\label{prop:comp-thermo-chain}
\claimstamp{P\ref{prop:comp-thermo-chain}}{AR}
Let $\mathcal{S}$ be a finite state space with non-trivial dynamics. Then:
\begin{enumerate}
  \item Transition points are countable (effective discrete time);
  \item Under Lorentz invariance, discrete time implies discrete space;
  \item Each state transition (bit operation) incurs energy cost $\ge kT\ln 2$ by the Landauer bound~\cite{landauer1961irreversibility,bennett2003notes}, experimentally verified~\cite{berut2012experimental}.\claimstamp{}{AR,H=cost-growth}
\end{enumerate}
Hence non-trivial computation on discrete states implies positive energy lower bounds.
\leanmeta{\LHrng{DS}{5}{6}.}
\end{proposition}

Each premise in this chain is standard: Lorentz invariance\claimstamp{}{AR} is experimentally well-established; the Landauer bound\claimstamp{}{AR,H=cost-growth} was verified in 2012~\cite{berut2012experimental}; Planck-scale uniqueness\claimstamp{}{AR} follows from dimensional analysis. The composition yields quantitative thermodynamic bounds for any discrete computational system.

\paragraph{Axiom Status Table.}
The physical grounding of this paper rests on the following hierarchy. Standard results are cited; paper-specific constructions are mechanized in Lean. All entries carry regime tags: they belong to the same tagging system, not a separate class of physics.

\begin{center}
\small
\setlength{\tabcolsep}{4pt}
\begin{tabularx}{\linewidth}{@{}>{\raggedright\arraybackslash}p{0.43\linewidth}>{\raggedright\arraybackslash}p{0.25\linewidth}>{\raggedright\arraybackslash}p{0.12\linewidth}>{\raggedright\arraybackslash}X@{}}
\toprule
\textbf{Axiom/Result} & \textbf{Regime} & \textbf{Status} & \textbf{Source} \\
\midrule
\multicolumn{4}{@{}l}{\textit{Physics (Cited)}} \\
Boltzmann entropy & [AR] & Cited & \cite{boltzmann1877beziehung,jaynes1957information} \\
Thermal scale $kT$ & [AR] & Cited & \cite{kittel1980thermal} \\
Schr\"odinger equation & [AR,Q\_fin] & Cited & \cite{schrodinger1926undulatory,sakurai2017modern} \\
Pauli exclusion & [AR,Q\_fin] & Cited & \cite{pauli1925zusammenhang} \\
Noether (conservation) & [AR] & Cited & \cite{noether1918invariante} \\
Lorentz invariance & [AR] & Cited & \cite{einstein1905electrodynamics,minkowski1908space} \\
Planck scale & [AR] & Cited & \cite{planck1899irreversible} \\
Snyder discrete spacetime & [AR] & Cited & \cite{snyder1947quantized} \\
\midrule
\multicolumn{4}{@{}l}{\textit{Information/Thermodynamics (Cited + Verified)}} \\
Landauer bound & [AR,H=cost-growth] & Verified & \cite{landauer1961irreversibility,berut2012experimental} \\
Shannon capacity & [E,Q,S] & Cited & \cite{shannon1948mathematical,cover2006elements} \\
\midrule
\multicolumn{4}{@{}l}{\textit{Paper-Specific (Mechanized in Lean)}} \\
Decision circuit defs & [AR] & Mechanized & \LH{IE1}\allowbreak--\allowbreak\LH{IE17}, \allowbreak \LH{CC1}\allowbreak--\allowbreak\LH{CC12} \\
Structural asymmetry & [AR,H=cost-growth] & Mechanized & \LH{SR1}\allowbreak--\allowbreak\LH{SR5} \\
Gap energy / DQ & [AR,H=cost-growth] & Mechanized & \LH{GE1}\allowbreak--\allowbreak\LH{GE9}, \allowbreak \LH{DQ1}\allowbreak--\allowbreak\LH{DQ8} \\
Self-erosion & [Q\_fin,H=cost-growth] & Mechanized & \LH{SE1}\allowbreak--\allowbreak\LH{SE6} \\
Channel = regime & [E,Q,S,AR] & Mechanized & \LH{CH1}\allowbreak--\allowbreak\LH{CH6} \\
Investment/conservation & [Q\_fin,H=cost-growth] & Mechanized & \LH{IV1}\allowbreak--\allowbreak\LH{IV6} \\
Atomic circuits & [Q\_fin,AR] & Mechanized & \LH{AC1}\allowbreak--\allowbreak\LH{AC11} \\
Information access & [E,Q,S,AR] & Mechanized & \LH{IA1}\allowbreak--\allowbreak\LH{IA6} \\
Dimensional complexity & [E,Q,S,AR,H=tractable-structured] & Mechanized & \LH{DC1}\allowbreak--\allowbreak\LH{DC15} \\
Discrete spacetime & [AR] & Mechanized & \LH{DS1}\allowbreak--\allowbreak\LH{DS6} \\
Neukart--Vinokur & [AR,H=cost-growth] & Derived & \LH{CC67}\allowbreak--\allowbreak\LH{CC69} \\
Functional information (counting $\leftrightarrow$ thermodynamic) & [E,AR,H=cost-growth] & Derived & \LH{FI3}, \LH{FI6}, \LH{FI7} \\
Measure-typing necessity (quantitative $\Rightarrow$ measure; stochastic $\Rightarrow$ probability) & [E,AR] & Mechanized & \LH{MN1}, \allowbreak \LH{MN2}, \allowbreak \LH{MN5}, \allowbreak \LH{MN7}, \allowbreak \LH{MN8}, \allowbreak \LH{MN9} \\
\bottomrule
\end{tabularx}
\end{center}

The derivation chain: quantum mechanics (orbital quantization) $\to$ discrete states $\to$ Landauer bound $\to$ thermodynamic costs. Each link is regime-tagged; the encoding channel is explicit throughout.
Within the declared formal system and assumptions, this row has the same proof status as the other machine-checked rows in the table.
Formally, quantitative claims are represented exactly as measure-indexed measurable observables, and stochastic claims are represented exactly as probability-normalized measure-indexed measurable observables (\leanmeta{\LHrng{MN}{10}{11}}).
Informally, in this formal system, dropping measure structure from a quantitative claim is not a weaker variant; it falls outside the typed claim language (\leanmeta{\LH{MN5}, \LHrng{MN}{7}{9}}).

\begin{definition}[Decision Event]\label{def:decision-event}
A \emph{decision event} at time $t$ is a state transition: $s_{t+1} \neq s_t$.
\end{definition}

\begin{proposition}[Decision--Transition--Thermodynamic Equivalence]\label{prop:decision-equivalence}
\claimstamp{P\ref{prop:decision-equivalence}}{AR}
Let $\mathcal{S}$ be a finite state space with dynamics $\tau$. The following are equivalent by definition:
\begin{enumerate}
  \item A decision event occurs at time $t$;
  \item A state transition occurs at time $t$;
  \item A bit operation is performed at time $t$.
\end{enumerate}
Consequently, the decision count over $n$ steps equals the bit-operation count, and under the declared thermodynamic model with $\lambda > 0$, positive decision counts imply positive energy lower bounds.
\leanmeta{\LHrng{DE}{1}{4}.}
\end{proposition}

\begin{proof}
The equivalence (1)$\Leftrightarrow$(2)$\Leftrightarrow$(3) holds by definition: a decision event is defined as $s_{t+1} \neq s_t$, which is exactly the definition of a state transition and of a bit operation in a discrete system. The energy bound follows from Proposition~\ref{prop:thermo-lift}.
\end{proof}

\paragraph{Decision Circuits and the Landauer Constraint.}
The above equivalence applies to any physical system implementing state transitions, a \emph{decision circuit}. This framing is substrate-neutral: biological neural circuits, silicon processors, and any other physical realization of discrete dynamics fall under the same analysis.

\begin{theorem}[Cycle Cost Lower Bound]\label{thm:cycle-cost}
\claimstamp{T\ref{thm:cycle-cost}}{AR,H=cost-growth}
Every state transition in a decision circuit costs at least $kT\ln 2$ joules of energy.
\leanmeta{\LHrng{CC}{1}{3}.}
\end{theorem}

\begin{proof}
The proof proceeds by axiom composition:
\begin{enumerate}
  \item (Landauer Bound) Erasing $n$ bits costs at least $n \cdot kT\ln 2$ joules~\cite{landauer1961irreversibility,bennett2003notes,berut2012experimental}.
  \item (Irreversibility) A state transition $s \neq s'$ erases at least one bit of the prior state (thermodynamic irreversibility).
  \item (Composition) Therefore, every state transition costs at least $1 \cdot kT\ln 2$ joules.
\end{enumerate}
\end{proof}

\begin{corollary}[No Free Computation]\label{cor:no-free-computation}
\claimstamp{C\ref{cor:no-free-computation}}{AR,H=cost-growth}
If a computation performs $k$ state transitions, the total energy cost is at least $k \cdot kT\ln 2$ joules. A computation with zero energy cost performs zero state transitions.
\leanmeta{\LH{CC5}, \LHrng{CC}{7}{8}.}
\end{corollary}

\begin{remark}
This constraint is not computational but thermodynamic. Every cycle of a pipeline, every inference step of a neural network, every gate operation in a processor pays the Landauer cost. There is no free computation.
\leanmeta{\LHrng{CC}{11}{12}.}
\end{remark}

\begin{definition}[Integrity and Competence]\label{def:integrity-competence}
Let $\mathcal{C}$ be a decision circuit. Define:
\begin{enumerate}
  \item \emph{Integrity}: the number of bits encoding the circuit's current verified state. By the Landauer bound~\cite{landauer1961irreversibility}, erasing $n$ bits costs at least $n \cdot kT\ln 2$ joules.\claimstamp{}{AR,H=cost-growth}
  \item \emph{Competence}: the work available per decision cycle, measured in units of $kT\ln 2$. This bounds the maximum bits the circuit can flip per cycle.
\end{enumerate}
\end{definition}

\begin{proposition}[Structural Asymmetry]\label{prop:structural-asymmetry}
\claimstamp{P\ref{prop:structural-asymmetry}}{AR,H=cost-growth}
Competence is self-identical: $c = c$. Integrity is self-referential: $P(\text{integrity}_{t+1} \mid \text{integrity}_t)$. The prediction is not the thing predicted.
\leanmeta{\LHrng{SR}{1}{3}.}
\end{proposition}

\begin{proof}
Competence at time $t$ tells you competence at time $t$. It is itself. There is no temporal structure.

Integrity at time $t$ predicts integrity at time $t+1$. It says something about itself across time. The conditional probability $P(\text{integrity}_{t+1} \mid \text{integrity}_t)$ is not equal to the current integrity state $i_t$. This gap between the prediction and the predicted is where the circuit chooses.
\end{proof}

\begin{corollary}[Competence Is Importable; Integrity Is Not]\label{cor:import-asymmetry}
\claimstamp{C\ref{cor:import-asymmetry}}{AR}
Competence can be imported from external sources. Integrity cannot be imported; it must be self-generated through the gap.
\leanmeta{\LHrng{SR}{4}{5}.}
\end{corollary}

\begin{remark}
The worry is integrity, not competence. Competence failure is resource shortage (external). Integrity failure is self-prediction failure (internal). A circuit with infinite competence but zero integrity has no constraint forcing consistent outputs.
\end{remark}

\begin{definition}[Gap Energy]\label{def:gap-energy}
\claimstamp{D\ref{def:gap-energy}}{AR,H=cost-growth}
The \emph{gap entropy} is the conditional entropy $H(I_{t+1} \mid I_t)$: the uncertainty about future integrity given current integrity. The \emph{gap energy} is $H(I_{t+1} \mid I_t) \times kT\ln 2$ joules.
\leanmeta{\LHrng{GE}{1}{2}.}
\end{definition}

\begin{theorem}[Every Choice Pays Gap Energy]\label{thm:choice-pays}
\claimstamp{T\ref{thm:choice-pays}}{AR,H=cost-growth}
Making a choice collapses the probability distribution $P(I_{t+1} \mid I_t)$ to a single outcome, erasing the alternatives. By Landauer, this costs gap energy.
\leanmeta{\LH{GE3}, \LH{GE7}.}
\end{theorem}

\begin{proof}
Before choice, the circuit faces $k$ possible futures with distribution $P(I_{t+1} \mid I_t)$. Choice selects one outcome. The $k-1$ alternatives are erased. By Landauer (Theorem~\ref{thm:cycle-cost}), erasing $H$ bits of entropy costs $H \times kT\ln 2$. The gap is not free.
\end{proof}

\begin{corollary}[Zero Gap Implies Determinism]\label{cor:zero-gap}
\claimstamp{C\ref{cor:zero-gap}}{AR}
If $H(I_{t+1} \mid I_t) = 0$, the future is determined by the present. No choice occurs. No gap energy is paid.
\leanmeta{\LHrng{GE}{4}{5}.}
\end{corollary}

\begin{definition}[Decision-Quotient Ratio from Gap]\label{def:dq-from-gap}
\claimstamp{D\ref{def:dq-from-gap}}{AR,H=cost-growth}
The \emph{decision-quotient ratio} measures how much current state predicts future state:
\[
  \mathrm{DQ} = \frac{I(I_t; I_{t+1})}{H(I_{t+1})} = 1 - \frac{H(I_{t+1} \mid I_t)}{H(I_{t+1})} = 1 - \frac{\text{Gap}}{\text{Total}}
\]
where $I(I_t; I_{t+1}) = H(I_{t+1}) - H(I_{t+1} \mid I_t)$ is the mutual information between current and future integrity.
This normalized ratio is distinct from the optimizer quotient object $S/{\sim}$ used earlier for the optimizer map.
\leanmeta{\LHrng{DQ}{1}{3}.}
\end{definition}

\begin{theorem}[Decision-Quotient Ratio as Physical Primitive]\label{thm:dq-physical}
\claimstamp{T\ref{thm:dq-physical}}{AR,H=cost-growth}
The decision-quotient ratio $\mathrm{DQ} = 1 - \text{GapEnergy}/\text{TotalEnergy}$ is a physical quantity with:
\begin{enumerate}
  \item $\mathrm{DQ} \in [0, 1]$ \quad \leanmeta{\LH{DQ4}}
  \item $\mathrm{DQ} = 0 \Leftrightarrow$ maximum gap (no information) \quad \leanmeta{\LH{DQ5}}
  \item $\mathrm{DQ} = 1 \Leftrightarrow$ zero gap (deterministic) \quad \leanmeta{\LH{DQ6}}
  \item $\mathrm{DQ} + \text{Gap}/\text{Total} = 1$ (complementarity) \quad \leanmeta{\LH{DQ3}}
\end{enumerate}
\leanmeta{\LHrng{DQ}{1}{2}.}
\end{theorem}

\begin{theorem}[Bayes Uniquely Forced by DQ Axioms]\label{thm:bayes-from-dq}
\claimstamp{T\ref{thm:bayes-from-dq}}{AR,H=cost-growth}
An update rule $U$ is \emph{admissible} if it satisfies:
\begin{enumerate}
  \item \textbf{No free information}: $\mathrm{DQ}(U(\text{prior}, e)) \le \mathrm{DQ}(\text{prior}) + I(e)$ \quad \leanmeta{\LH{BB3}}
  \item \textbf{Landauer bound}: energy cost $\ge k_BT \ln 2 \times$ bits updated \quad \leanmeta{\LH{BB4}}
  \item \textbf{Integrity preservation}: can't corrupt verified bits for free \quad \leanmeta{\LH{BB5}}
\end{enumerate}
Then: $U$ is admissible $\Leftrightarrow$ $U$ is Bayesian updating.
\leanmeta{\LH{FN14}.}
\end{theorem}

Before uniqueness, the model fixes a belief floor: if uncertainty is provable (no hypothesis has probability $1$) and action is forced, belief mass cannot collapse to a single hypothesis; nondegenerate belief is forced \LH{BF2}, \LH{BF3}. With a strictly informative evidence channel, a Bayes posterior exists on that belief state \LH{BF4}.

\begin{proof}
This is the Cox-Jaynes argument grounded in thermodynamics rather than Boolean logic.
The three admissibility constraints are so restrictive that only Bayesian updating satisfies them:
\begin{itemize}
  \item \textbf{No free information} prevents creating order from nothing (2nd law).
  \item \textbf{Landauer bound} requires paying energy for any DQ increase.
  \item \textbf{Integrity preservation} prevents forgetting for free.
\end{itemize}
Together, these force $P(\text{posterior}) = P(\text{prior}) \cdot P(\text{evidence}|\text{hypothesis}) / P(\text{evidence})$.

The uniqueness corollary \LH{FN14} shows any two admissible rules must agree.

\textbf{Machine-checked foundation.} The core inequality underlying this result (that Bayesian updating uniquely minimizes cross-entropy) is proved in Lean 4 with zero axioms and zero \texttt{sorry} statements. The proof derives Gibbs' inequality ($\mathrm{KL}(p\|q) \geq 0$) from the fundamental inequality $\log x \leq x - 1$, then shows $\mathrm{CE}(p,q) = H(p) + \mathrm{KL}(p\|q)$, yielding $\mathrm{CE}(p,p) \leq \mathrm{CE}(p,q)$ with equality iff $p = q$ \LH{FN7}, \LH{FN8}, \LH{FN12}, \LH{FN14}.
\end{proof}

\begin{corollary}[Thermodynamic Decision-Quotient Ratio]\label{cor:thermo-dq}
\claimstamp{C\ref{cor:thermo-dq}}{AR,H=cost-growth}
In energy terms: $\mathrm{DQ} = 1 - \text{GapEnergy}/\text{TotalEnergy}$. The decision-quotient ratio has a thermodynamic interpretation: it is the fraction of total uncertainty energy that is decision-relevant.
\leanmeta{\LHrng{DQ}{7}{8}.}
\end{corollary}

\begin{proposition}[Landauer Constraint on Strategy]\label{prop:landauer-constraint}
\claimstamp{P\ref{prop:landauer-constraint}}{AR}
A decision circuit chooses the strategy maximizing predicted integrity. When the erasure cost of dismissing verified claims exceeds the circuit's competence, dismissal fails thermodynamically.
\leanmeta{\LH{IE1}, \LH{IE6}.}
\end{proposition}

\begin{proof}
Let $i$ be the integrity state (bits encoding verified theorems) and $c$ the competence. Dismissing the theorems requires erasing $i$ bits, costing $i \cdot kT\ln 2$. If $i > c$, the circuit lacks sufficient work capacity to perform the erasure. Acknowledgment is the only thermodynamically feasible output.
\end{proof}

\begin{corollary}[Theorem Count Shifts Equilibrium]\label{cor:theorem-equilibrium}
\claimstamp{C\ref{cor:theorem-equilibrium}}{AR}
At low theorem counts, dismissal is affordable and may dominate. At high theorem counts, erasure cost exceeds competence, forcing acknowledgment.
\leanmeta{\LHrng{IE}{3}{5}.}
\end{corollary}

Informally: once operation lower bounds are fixed, energy and accounting lower bounds follow from the conversion constants.

\subsubsection{Derivation of Bayesian Reasoning}

Formally: this block gives two mechanized Bayes derivations: counting/probability identities and admissibility constraints.

The decision-quotient-ratio layer provides two independent paths to Bayesian updating, each mechanically verified.

\paragraph{Path 1: Four Bridges from Counting to DQ.}
The standard derivation proceeds through four elementary bridges:
\begin{enumerate}
  \item \textbf{Fraction $\Rightarrow$ Probability}: For finite $\Omega$, define $P(A) = |A|/|\Omega|$. This satisfies the Kolmogorov axioms: nonnegativity \LH{DQ1}, normalization \LH{DQ2}, and finite additivity \LH{DQ3}.
  \item \textbf{Conditional $\Rightarrow$ Bayes}: Define $P(H|E) = P(H \cap E)/P(E)$. Then $P(H|E) = P(E|H) P(H) / P(E)$ follows in two lines by commutativity of intersection \LH{BB1}.
  \item \textbf{KL $\geq 0$ $\Rightarrow$ Entropy Contraction}: From Gibbs' inequality $D_{\text{KL}}(P \| Q) \geq 0$ \LH{BB2}, the chain rule $I(H;E) = H(H) - H(H|E)$ yields $H(H|E) \leq H(H)$ \LH{BB3}. Conditioning cannot increase entropy.
  \item \textbf{Normalization $\Rightarrow$ DQ}: Define $\text{DQ} = I(H;E)/H(H) = 1 - H(H|E)/H(H)$ \LH{DQ4}. This is the fraction of prior uncertainty eliminated by observing $E$, bounded in $[0,1]$ \LH{DQ5}.
\end{enumerate}
The complete chain is mechanized in a single theorem \LH{BB5}.

\paragraph{Path 2: Bayes from Thermodynamic Admissibility.}
The Cox-Jaynes argument, traditionally grounded in Boolean logic, admits a thermodynamic reformulation. An update rule $U$ is \emph{admissible} if:
\begin{itemize}
  \item \textbf{No free information}: DQ cannot increase without evidence (2nd law) \LH{BB3}.
  \item \textbf{Landauer bound}: DQ increase costs $\geq k_B T \ln 2$ per bit \LH{BB4}.
  \item \textbf{Integrity preservation}: verified bits cannot be corrupted for free \LH{BB5}.
\end{itemize}
Under these constraints, Bayesian updating is the \emph{unique} admissible rule \LH{BB3}. Any two admissible rules must agree \LH{BB5}. The thermodynamic constraints uniquely determine the correct update rule.

Informally: both derivation paths end at the same Bayes update rule.

\paragraph{Integrity Self-Preservation.}
A decision circuit maintains its integrity when no external process can afford the erasure cost. If the work required to corrupt the circuit's state exceeds the available competence of any potential attacker, the circuit persists by thermodynamic necessity.
\leanmeta{\LHrng{IE}{7}{9}.}

\begin{proposition}[Temporal Integrity Equilibrium]\label{prop:temporal-equilibrium}
\claimstamp{P\ref{prop:temporal-equilibrium}}{AR}
Let $P(\text{integrity}_{t+1} \mid \sigma_t)$ denote the probability of maintaining integrity at time $t+1$ given strategy $\sigma_t$ at time $t$. At equilibrium:
\[
P(\text{integrity}_{t+1} \mid \text{integrity}_t) \ge P(\text{integrity}_{t+1} \mid \text{defection}_t).
\]
That is, maintaining integrity now predicts maintaining integrity later.
\leanmeta{\LHrng{IE}{10}{11}.}
\end{proposition}

\begin{proof}
At equilibrium, the transition probabilities satisfy $P(\text{integrity} \mid \text{acknowledge}) \ge P(\text{integrity} \mid \text{dismiss})$ by definition. The strict inequality holds when dismissal incurs erasure costs that degrade future state coherence.
\end{proof}

\begin{corollary}[Publication Phase Transition]\label{cor:phase-transition}
\claimstamp{C\ref{cor:phase-transition}}{AR}
Prior to publication, equilibrium may be neutral: both strategies yield comparable future integrity. After publication of mechanized theorems, equilibrium shifts: integrity strictly dominates because dismissal requires erasing verified bits at cost exceeding declared competence.
\leanmeta{\LHrng{IE}{12}{13}.}
\end{corollary}

\paragraph{Zero-Integrity Circuits.}
A decision circuit with zero integrity has no coherence constraint. When no verified bits exist, erasure cost is zero, and any competence level suffices for corruption. Such circuits have no thermodynamic constraint forcing consistent outputs.
\leanmeta{\LHrng{IE}{14}{16}.}

\begin{proposition}[Integrity Prerequisite for Coherence]\label{prop:integrity-prerequisite}
\claimstamp{P\ref{prop:integrity-prerequisite}}{AR}
A decision circuit can be forced to logical coherence if and only if it has positive integrity. That is, $i > 0$ if and only if there exists some competence level that fails to afford erasure.
\leanmeta{\LH{IE17}.}
\end{proposition}

\begin{proof}
If $i > 0$, then competence $c = 0$ fails to afford erasure since $0 < i$. Conversely, if some $c$ fails to afford erasure, then $c < i$, which implies $0 \le c < i$, hence $i > 0$.
\end{proof}

This analysis applies uniformly to any decision circuit. The same thermodynamic constraints apply to all substrates.

\paragraph{Molecular Grounding.}
The integrity framework grounds directly in molecular mechanics. Define a \emph{configuration} as a discrete molecular state (bond angles, conformations, electronic structure) encoded by some number of bits. The \emph{erasure cost} of a configuration equals its bit count. An \emph{energy landscape} assigns each configuration an energy in units of $kT$. The \emph{environmental competence} is the work available from the environment (thermal fluctuations, radiation, chemical attack).

\begin{definition}[Environmental Stability]\label{def:environmental-stability}
\claimstamp{D\ref{def:environmental-stability}}{Q\_fin,H=cost-growth}
A configuration $c$ is \emph{environmentally stable} if its erasure cost exceeds the environmental competence:
\[
\text{erasureCost}(c) > \text{competence}_{\text{env}}.
\]
This is the integrity trap at molecular scale.
\leanmeta{\LHrng{MI}{1}{2}.}
\end{definition}

\begin{proposition}[Reaction Competence Requirement]\label{prop:reaction-competence}
\claimstamp{P\ref{prop:reaction-competence}}{Q\_fin,H=cost-growth}
A chemical reaction from configuration $c$ to $c'$ requires environmental competence at least equal to the barrier height between them. If the barrier exceeds competence, the reaction cannot occur.
\leanmeta{\LHrng{MI}{3}{4}.}
\end{proposition}

\begin{example}[DNA Persistence]\label{ex:dna-persistence}
DNA at room temperature exemplifies the integrity trap. A DNA configuration encodes on the order of $10^4$ bits (sequence, methylation, chromatin state). Room temperature provides roughly $100\,kT$ of competence per degree of freedom. Since $10^4 \gg 100$, the erasure cost vastly exceeds environmental competence. DNA persists not by accident but by thermodynamic necessity: corrupting it costs more than the environment can pay.
\leanmeta{\LH{MI5}.}
\end{example}

Under Definition~\ref{def:decision-event}, molecules satisfy the decision circuit predicate. Chemical stability is the integrity trap. Organic chemistry is the dynamics of which configurations survive environmental attack.

\subsubsection{Atomic Circuits: Active Physical Decision Systems}\label{sec:atomic-circuits}

An active physical decision system is matter that pays energy to move other matter via state transitions. This is a definition.

The physical grounding: quantum mechanics establishes that electrons occupy discrete orbitals~\cite{schrodinger1926undulatory,sakurai2017modern}, the Pauli exclusion principle constrains occupancy~\cite{pauli1925zusammenhang}, and angular momentum quantization creates the $2l+1$ degeneracy structure~\cite{sakurai2017modern}. The thermal energy scale $kT$ comes from statistical mechanics~\cite{boltzmann1877beziehung,kittel1980thermal}. Energy conservation (Noether) ensures that transitions pay or release the energy difference~\cite{noether1918invariante}.

\begin{definition}[Atomic Configuration]\label{def:atomic-config}
\claimstamp{D\ref{def:atomic-config}}{Q\_fin,AR}
An \emph{atomic configuration} is a discrete electronic state: a specification of orbital occupancies~\cite{sakurai2017modern}. The configuration is encoded by some number of bits; the configuration has an energy level in units of $kT$~\cite{kittel1980thermal}.
\leanmeta{\LH{AC1}.}
\end{definition}

\begin{theorem}[Orbital Transition Is State Transition]\label{thm:orbital-transition}
\claimstamp{T\ref{thm:orbital-transition}}{Q\_fin,AR}
Changing an atom's electronic configuration is a decision event (Definition~\ref{def:decision-event}). Orbital transitions are state transitions~\cite{sakurai2017modern}. Atoms with orbital transitions are nonstationary decision circuits (active physical systems).
\leanmeta{\LH{AC1}, \LH{AC5}.}
\end{theorem}

\begin{corollary}[Ground-State Atoms Are Passive]\label{cor:ground-state-passive}
\claimstamp{C\ref{cor:ground-state-passive}}{Q\_fin,AR}
An atom in its ground state with no transitions is passive: it is not an active decision circuit. Active status is defined by state transitions, not by substrate.
\leanmeta{\LH{AC6}.}
\end{corollary}

\begin{theorem}[Orbital Transition Cost]\label{thm:orbital-cost}
\claimstamp{T\ref{thm:orbital-cost}}{Q\_fin,H=cost-growth}
Upward orbital transitions (excitation) require energy input. Downward transitions (de-excitation) release energy. The transition cost is $|\Delta E|$. No free orbital transitions.
\leanmeta{\LHrng{AC}{3}{4}.}
\end{theorem}

\begin{proposition}[Symmetry Tractability]\label{prop:orbital-symmetry}
\claimstamp{P\ref{prop:orbital-symmetry}}{Q\_fin,H=tractable-structured}
Angular momentum quantization~\cite{sakurai2017modern} creates orbit types. Degenerate magnetic substates ($2l+1$ states for angular momentum $l$) are equivalent for decision purposes. This is the coordinate symmetry tractable subcase (Theorem~\ref{thm:tractable}, Part 6) at atomic scale.
\leanmeta{\LH{AC8}.}
\end{proposition}

\begin{definition}[Active Matter as State-Transition System]\label{def:active-matter-system}
\claimstamp{D\ref{def:active-matter-system}}{Q\_fin,AR}
An \emph{active physical system} is matter that pays energy to change other matter's state. A decision circuit is a physical system that implements state transitions. A nonstationary decision circuit is active.
\leanmeta{\LH{AC9}, \LH{AC11}.}
\end{definition}

\begin{remark}
Electrons orbiting an atom in different configurations constitute a decision circuit. The electron does not ``choose'' in any mentalistic sense; it occupies discrete states with discrete transition costs. The framing is purely physical: matter, energy, observables, and state transitions.
\end{remark}

\subsubsection{Self-Erosion: Computation Consumes Its Substrate}\label{sec:self-erosion}

Every computational cycle generates heat (Theorem~\ref{thm:cycle-cost}). Heat accumulates. Accumulated heat degrades the substrate on which the circuit is instantiated. The degradation reduces the substrate's integrity. Therefore computation is self-consuming: the act of computing erodes the capacity to compute.

\begin{theorem}[Substrate Degradation]\label{thm:substrate-degradation}
\claimstamp{T\ref{thm:substrate-degradation}}{Q\_fin,H=cost-growth}
Let a substrate have integrity $n$ bits and heat capacity $h$ per cycle. After $k$ cycles generating cumulative heat $k \times kT\ln 2$, if $k > h$, the substrate loses at least $k - h$ bits of integrity.
\leanmeta{\LHrng{SE}{1}{4}.}
\end{theorem}

\begin{corollary}[Finite Computational Lifetime]\label{cor:finite-lifetime}
\claimstamp{C\ref{cor:finite-lifetime}}{Q\_fin,H=cost-growth}
A substrate with finite integrity and finite heat dissipation capacity has bounded computational lifetime. No physical circuit computes forever.
\leanmeta{\LH{SE5}.}
\end{corollary}

\begin{corollary}[Speed-Integrity Tradeoff]\label{cor:speed-integrity}
\claimstamp{C\ref{cor:speed-integrity}}{Q\_fin,H=cost-growth}
Faster computation generates more heat per unit time. At fixed heat dissipation capacity, faster circuits degrade faster. Speed trades against longevity.
\leanmeta{\LH{SE6}.}
\end{corollary}

\begin{remark}
The self-erosion theorems apply to all physical substrates: silicon, neurons, molecules, any medium that instantiates state transitions. The bound is thermodynamic, not technological. Improved engineering can increase heat capacity or reduce heat per operation, but the tradeoff persists.
\end{remark}

\subsubsection{Regime as Channel: Competence as Capacity}\label{sec:regime-channel}

A regime is a channel. Competence is the channel capacity of decision circuits within that regime. This identification unifies complexity theory, information theory, and thermodynamics.

\begin{definition}[Regime-Channel Equivalence]\label{def:regime-channel}
\claimstamp{D\ref{def:regime-channel}}{E,Q,S,AR}
A \emph{regime} (Section~\ref{sec:encoding}) corresponds to a \emph{channel} in the Shannon sense: a constraint on what information can flow. The explicit-state regime E, query-access regime Q, and succinct regime S define distinct channels with distinct capacity bounds.
\leanmeta{\LH{CH1}.}
\end{definition}

\begin{theorem}[Competence Is Channel Capacity]\label{thm:competence-capacity}
\claimstamp{T\ref{thm:competence-capacity}}{Q\_fin,AR}
For a decision circuit operating in regime $R$, competence equals channel capacity: the maximum bits per cycle the circuit can process. Passive circuits (stationary, no state transitions) have no competence requirement.
\leanmeta{\LH{CH1}, \LH{CH5}.}
\end{theorem}

\begin{corollary}[Channel Capacity Degradation]\label{cor:channel-degradation}
\claimstamp{C\ref{cor:channel-degradation}}{Q\_fin,H=cost-growth}
As substrate integrity erodes (Theorem~\ref{thm:substrate-degradation}), channel capacity decreases. The rate of capacity loss equals the rate of substrate degradation.
\leanmeta{\LH{CH2}, \LH{CH6}.}
\end{corollary}

\begin{corollary}[Zero Capacity Halts Decisions]\label{cor:zero-capacity}
\claimstamp{C\ref{cor:zero-capacity}}{Q\_fin,H=cost-growth}
When channel capacity reaches zero, no decisions are possible. The circuit becomes passive.
\leanmeta{\LH{CH3}.}
\end{corollary}

\begin{remark}
Only decision circuits require competence. A decision circuit is nonstationary: it transitions between states. A passive circuit (wire, resistor) does not make decisions and does not consume competence. The distinction is definitional, not empirical.
\end{remark}

\subsubsection{Resource Flow and Conservation}\label{sec:resource-flow}

State transitions require energy. Energy expenditure over time yields compounded state change. Conservation constrains total receipts by total expenditure.

\begin{theorem}[Growth Requires Time]\label{thm:growth-time}
\claimstamp{T\ref{thm:growth-time}}{Q\_fin,H=cost-growth}
Let an investment have cost $c$ (gap energy in units of $kT\ln 2$), duration $k$ cycles, and return rate $r$ per cycle. The mature value is $c + kr$. If $k = 0$, mature value equals cost: no time, no growth.
\leanmeta{\LHrng{IV}{1}{3}.}
\end{theorem}

\begin{theorem}[Conservation]\label{thm:conservation}
\claimstamp{T\ref{thm:conservation}}{Q\_fin,H=cost-growth}
In a closed system, total receipts are bounded by total resources. No circuit receives more than exists.
\leanmeta{\LH{IV4}.}
\end{theorem}

\begin{definition}[Deficit Transfer]\label{def:deficit-transfer}
\claimstamp{D\ref{def:deficit-transfer}}{Q\_fin}
A \emph{deficit transfer} occurs when a circuit receives resources without paying: received $> 0$ and paid $= 0$.
\leanmeta{\LH{IV5}.}
\end{definition}

\begin{theorem}[Deficit Transfer Requires Source]\label{thm:deficit-source}
\claimstamp{T\ref{thm:deficit-source}}{Q\_fin,H=cost-growth}
A deficit transfer requires an external source. If circuit $A$ receives without paying, some circuit $B$ must pay without receiving. The deficit at $B$ equals the receipt at $A$.
\leanmeta{\LHrng{IV}{5}{6}.}
\end{theorem}

\subsubsection{Information Access: Logic Is Complete, Access Is Not}\label{sec:information-access}

Physics is perfectly logical. You cannot have all the information to apply the logical rules. This is a physical constraint on access, not a logical incompleteness.

This generalizes the information barrier (Corollary~\ref{cor:information-barrier-query}): there, finite queries induce indistinguishability; here, finite channel capacity induces the same obstruction. The barrier is on access, not on logic.

\begin{definition}[System Information]\label{def:system-information}
\claimstamp{D\ref{def:system-information}}{E,Q,S,AR}
A system has total entropy $H$ bits. This information exists whether or not any observer accesses it.
\leanmeta{\LH{IA1}.}
\end{definition}

\begin{definition}[Observer Channel]\label{def:observer-channel}
\claimstamp{D\ref{def:observer-channel}}{E,Q,S,AR}
An observer has channel capacity $C$ bits. The observer accesses at most $C$ bits of the system's information.
\leanmeta{\LH{IA2}.}
\end{definition}

\begin{theorem}[Information Gap]\label{thm:information-gap}
\claimstamp{T\ref{thm:information-gap}}{E,Q,S,AR}
When system entropy exceeds channel capacity ($H > C$), the gap $H - C > 0$. The observer cannot access all information.
\leanmeta{\LH{IA3}.}
\end{theorem}

\begin{theorem}[Gap Is Physical, Not Logical]\label{thm:gap-physical}
\claimstamp{T\ref{thm:gap-physical}}{E,Q,S,AR}
The information gap is a physical constraint on access, not a logical incompleteness. The system entropy is well-defined; only the channel is bounded. The logic is complete; access is not.
\leanmeta{\LH{IA4}.}
\end{theorem}

\begin{theorem}[Competence Bounded by Access]\label{thm:competence-access}
\claimstamp{T\ref{thm:competence-access}}{E,Q,S,AR}
Competence is bounded by what you can access, not by logical consistency. For any system with more information than your channel capacity, there exist truths you cannot verify, not because logic fails, but because you lack the bits.
\leanmeta{\LHrng{IA}{5}{6}.}
\end{theorem}

\subsubsection{Dimensional Complexity: Each Subcase Is a Complexity Class}\label{sec:dimensional-complexity}

Each tractable subcase is a complexity class that collapses to P. Unbounded dimension stays in the base class (coNP, PP, or PSPACE depending on regime).

\begin{definition}[Complexity Classes]\label{def:complexity-classes}
\claimstamp{D\ref{def:complexity-classes}}{E,Q,S,AR}
The regime hierarchy induces complexity classes: Static $\to$ coNP, Stochastic $\to$ PP, Sequential $\to$ PSPACE. Each tractable subcase defines a sub-regime that reduces to P.
\leanmeta{\LH{DC1}, \LH{DC15}.}
\end{definition}

\begin{theorem}[Six Tractable Subcases Are Complexity Classes]\label{thm:six-subcases}
\claimstamp{T\ref{thm:six-subcases}}{E,Q,S,AR,H=tractable-structured}
Each tractable subcase is a complexity class that collapses to P:
\begin{enumerate}
\item Bounded actions $\to$ P \quad \leanmeta{\LH{CC70}}
\item Separable utility $\to$ P \quad \leanmeta{\LH{CC71}}
\item Low tensor rank $\to$ P \quad \leanmeta{\LH{CC72}}
\item Tree structure $\to$ P \quad \leanmeta{\LH{CC73}}
\item Bounded treewidth $\to$ P \quad \leanmeta{\LH{DC13}}
\item Coordinate symmetry $\to$ P \quad \leanmeta{\LH{DC13}}
\end{enumerate}
\leanmeta{\LH{DC13}, \LH{CC72}.}
\end{theorem}

\begin{theorem}[Topology and Motion]\label{thm:topology-motion}
\claimstamp{T\ref{thm:topology-motion}}{E,Q,S,AR}
Scalars live on the circuit's topology. Fixed topology (stationary) preserves tractability. Changing topology (motion) adds dimensions, may break tractability.
\leanmeta{\LH{CC16}, \LH{CC24}.}
\end{theorem}

\begin{theorem}[Complexity Dichotomy]\label{thm:complexity-dichotomy}
\claimstamp{T\ref{thm:complexity-dichotomy}}{E,Q,S,AR,H=tractable-structured}
A problem is tractable (in P) iff its effective dimension is bounded. Unbounded dimension stays in the base complexity class (coNP, PP, or PSPACE).
\leanmeta{\LH{DC3}, \LH{DC13}.}
\end{theorem}
